
\section{Method}

\subsection{Overview}

This section describes our reproduction of MLA from DeepSeek-AI (2024). We 
do not introduce novel components; the goal is to (a) faithfully implement 
the original method and (b) generalize its weight-transfer to 
non-DeepSeek architectures.

\textbf{Figure 4.1} shows the MLA pipeline.

\begin{lstlisting}
   +----------------------------------------------+
   |  Input hidden h\_t $\in$ R^d                      |
   +--------------------+-------------------------+
                        |
        +---------------+---------------+
        |                               |
                                         
   +----------+                    +----------+
   |  Q proj  |                    |  kv\_down |  $\leftarrow$ latent down-projection
   |  W^Q     |                    |  W^DKV   |
   +----+-----+                    +----+-----+
        |                               |
        |                                
        |                          c\_t $\in$ R^c
        |                               |
        |                    +----------+----------+
        |                    |                     |
        |                                           
        |              +----------+          +----------+
        |              |  k\_up    |          |  v\_up    |
        |              |  W^UK    |          |  W^UV    |
        |              +----+-----+          +----+-----+
        |                   |                     |
        +-----------+-------+---------------------+
                    |
                     
              Attention(Q, K, V)
                    |
                     
              +----------+
              |  out\_proj|
              |  W^O     |
              +----+-----+
                   |
                    
              Output Y
\end{lstlisting}

\subsection{Forward Pass}

Given hidden states $h \in \mathbb{R}^{B \times L \times d}$:

\[
\begin{aligned}
q &= W^Q h \in \mathbb{R}^{B \times L \times H d\_h} \\
c &= W^{DKV} h \in \mathbb{R}^{B \times L \times c} \\
k &= W^{UK} c \in \mathbb{R}^{B \times L \times H\_{kv} d\_h} \\
v &= W^{UV} c \in \mathbb{R}^{B \times L \times H\_{kv} d\_h}
\end{aligned}
\tag{4.1}
\]

Then standard scaled dot-product attention with optional GQA repeat and 
RoPE.

\subsection{SVD-Based Initialization}

The core of our reproduction is the SVD-based initialization of 
$W^{DKV}, W^{UK}, W^{UV}$ from the concatenated pretrained $W^K, W^V$.

\subsubsection{Concatenation}

Given pretrained $W^K, W^V \in \mathbb{R}^{H_{kv} d_h \times d}$:

\[
W\_\text{fused} = \begin{pmatrix} W^K \\ W^V \end{pmatrix} \in \mathbb{R}^{2 H\_{kv} d\_h \times d}
\tag{4.2}
\]

\subsubsection{Truncated SVD}

\[
W\_\text{fused} \approx U \Sigma V^\top
\tag{4.3}
\]

where $U \in \mathbb{R}^{2 H_{kv} d_h \times r}$, $\Sigma \in \mathbb{R}^{r \times r}$, 
$V^\top \in \mathbb{R}^{r \times d}$, and $r = \min(c, 2 H_{kv} d_h, d)$.

\subsubsection{Setting the Parameters}

\[
W^{DKV} = V^\top, \qquad
W^{UK} = U\_{:H\_{kv} d\_h, :} \Sigma, \qquad
W^{UV} = U\_{H\_{kv} d\_h:, :} \Sigma
\tag{4.4}
\]

\subsubsection{Singular Value Normalization}

To preserve the total reconstruction scale, we compute:

\[
\gamma = \frac{\|W\_\text{fused}\|\_F}{\|\Sigma\|\_F}
\tag{4.5}
\]

and multiply both up-projections by $\gamma$:

\[
W^{UK} \leftarrow \gamma \cdot U\_{:H\_{kv} d\_h, :} \Sigma, \qquad
W^{UV} \leftarrow \gamma \cdot U\_{H\_{kv} d\_h:, :} \Sigma
\tag{4.6}
\]

This ensures:

\[
\|W^{UK} W^{DKV}\|\_F \approx \|W^K\|\_F, \qquad
\|W^{UV} W^{DKV}\|\_F \approx \|W^V\|\_F
\tag{4.7}
\]

\subsection{Weight Transfer from Pretrained Attention}

\subsubsection{Architecture-Aware Extraction}

We extract $(W^Q, W^K, W^V, W^O)$ differently per architecture:

\begin{table}[H]
\centering
\small
\begin{tabular}{ll}
\toprule
\textbf{Architecture} & \textbf{Extraction path} \\
\midrule
DistilBERT & \texttt{q\_lin, k\_lin, v\_lin, out\_lin} \\
RoBERTa & \texttt{self.query, self.key, self.value, output.dense} \\
GPT-2 & Split \texttt{c\_attn}; \texttt{c\_proj} for output \\
Qwen3.5 & \texttt{q\_proj, k\_proj, v\_proj, o\_proj} \\
\bottomrule
\end{tabular}
\end{table}

\subsubsection{Q and O Projections}

$W^Q$ and $W^O$ are copied directly (with LoRA adapters frozen at 
initialization):

\begin{lstlisting}
self.q\_proj.weight.data[:q\_dim, :hidden].copy\_(w\_q[:q\_dim, :hidden])
self.out\_proj.weight.data[:hidden, :q\_dim].copy\_(w\_o[:hidden, :q\_dim])
\end{lstlisting}

\subsubsection{K/V Concatenation and SVD}

\begin{lstlisting}
w\_kv = torch.cat([w\_k\_f[:kv\_dim, :hidden], w\_v\_f[:kv\_dim, :hidden]], dim=0)
u\_kv, s\_kv, vh\_kv = torch.linalg.svd(w\_kv, full\_matrices=False)
r = min(self.rank, s\_kv.shape[0], kv\_down.weight.shape[0], ...)

orig\_norm = torch.linalg.norm(w\_kv)
approx\_norm = torch.linalg.norm(s\_kv[:r])
norm\_scale = orig\_norm / max(approx\_norm, 1e-12)

self.kv\_down.weight.data[:r, :].copy\_(vh\_kv[:r, :])
up = (u\_kv[:, :r] * s\_kv[:r].unsqueeze(0)) * norm\_scale
self.k\_up.weight.data[:kv\_dim, :r].copy\_(up[:kv\_dim, :r])
self.v\_up.weight.data[:kv\_dim, :r].copy\_(up[kv\_dim:, :r])
\end{lstlisting}

\subsubsection{GQA Dimension Matching}

When $H_{kv} < H$ in the pretrained model, K and V have larger output 
dimension. We mean-pool before SVD:

\begin{lstlisting}
if w\_k\_f.shape[0] \textgreater{} self.kv\_dim:
    rep = w\_k\_f.shape[0] // self.kv\_dim
    w\_k\_f = w\_k\_f.view(self.kv\_dim, rep, -1).mean(dim=1)
\end{lstlisting}

\subsubsection{GPT-2 Conv1D}

Handled by transposing before SVD:

\begin{lstlisting}
if hasattr(src, "c\_attn"):
    if weight.shape[1] == 3 * self.hidden\_size:
        weight = weight.t()
\end{lstlisting}

\subsection{Nested Encoder Bypass}

For RoBERTa, the output projection is outside the inner attention. When 
nested, we set \texttt{is\_nested\_encoder = True} and skip \texttt{out\_proj} in forward.

When initialization is nested but $W^O$ is not extracted, we set 
\texttt{out\_proj} to an identity matrix:

\begin{lstlisting}
if self.is\_nested\_encoder:
    min\_d = min(self.out\_proj.weight.shape[0], self.out\_proj.weight.shape[1])
    self.out\_proj.weight.data.zero\_()
    self.out\_proj.weight.data[:min\_d, :min\_d].fill\_diagonal\_(1.0)
\end{lstlisting}

\subsection{Output Scale Calibration}

We introduce a trainable scalar \texttt{output\_scale} initialized to 1.0:

\begin{lstlisting}
self.output\_scale = nn.Parameter(torch.ones(1))
\end{lstlisting}

The forward pass multiplies the output by this scalar. This is \textbf{not} in 
the original DeepSeek-V2 paper, but absorbs distributional shift from the 
K/V factorization.

\subsection{LoRA Recovery}

We attach LoRA adapters to the Q and O projections:

\begin{lstlisting}
self.q\_proj = LoRALinear(..., lora\_rank=16, ...)
self.out\_proj = LoRALinear(..., lora\_rank=16, ...)
\end{lstlisting}

Only LoRA parameters (and \texttt{output\_scale}) are trainable. Base weights 
($W^{DKV}$, $W^{UK}$, $W^{UV}$, $W^Q$, $W^O$) are frozen.

\subsection{The Full Algorithm}

\textbf{Algorithm 1:} MLA initialization from pretrained attention
\begin{lstlisting}
Input:  pretrained attention module src
        latent dimension c
Output: MLA module with initialized weights

1:  Extract W\_Q, W\_K, W\_V, W\_O from src
2:  Orient to (out, in) form
3:  Handle Conv1D transposition
4:  GQA dimension match for K, V:
5:      if W\_K.shape[0] \textgreater{} H\_kv * d\_h: mean-pool
6:
7:  # Q and O direct copy
8:  q\_proj.weight[:q\_dim, :hidden] $\leftarrow$ W\_Q
9:  out\_proj.weight[:hidden, :q\_dim] $\leftarrow$ W\_O
10:
11: # K, V concatenate and SVD
12: W\_fused $\leftarrow$ concat([W\_K, W\_V], dim=0)
13: U, $\Sigma$, V^T $\leftarrow$ svd(W\_fused)
14: r $\leftarrow$ min(c, len($\Sigma$), ...)
15: $\gamma$ $\leftarrow$  W\_fused \_F /  $\Sigma$[:r] \_F
16:
17: kv\_down.weight[:r, :] $\leftarrow$ V^T[:r, :]
18: up $\leftarrow$ U[:, :r]   $\Sigma$[:r]   $\gamma$
19: k\_up.weight[:kv\_dim, :r] $\leftarrow$ up[:kv\_dim, :r]
20: v\_up.weight[:kv\_dim, :r] $\leftarrow$ up[kv\_dim:, :r]
21:
22: Freeze all base weights
23: Initialize LoRA A/B (B = 0)
24: return MLA module
\end{lstlisting}

\subsection{Complexity Analysis}

\begin{table}[H]
\centering
\small
\begin{tabular}{lll}
\toprule
\textbf{Phase} & \textbf{Time} & \textbf{Memory} \\
\midrule
Extraction & $\mathcal{O}(N_l d^2)$ & $\mathcal{O}(d^2)$ \\
SVD per layer & $\mathcal{O}((H d_h)^2 d)$ & $\mathcal{O}((H d_h) d)$ \\
Initialization & $\mathcal{O}(N_l c d)$ & $\mathcal{O}(c d)$ \\
Total & $\mathcal{O}(N_l (H d_h)^2 d)$ & $\mathcal{O}(N_l c d)$ \\
\bottomrule
\end{tabular}
\end{table}

\subsection{KV-Cache Reduction}

\begin{table}[H]
\centering
\small
\begin{tabular}{llll}
\toprule
\textbf{Model} & \textbf{Standard KV} & \textbf{MLA KV} & \textbf{Reduction} \\
\midrule
DistilBERT & $2 \cdot 12 \cdot 64 = 1536$ & 0.85 & \textbf{2.1$\times$} \\
RoBERTa & 1536 & 0.85 & \textbf{2.1$\times$} \\
GPT-2 & 1536 & 0.85 & \textbf{2.1$\times$} \\
\bottomrule
\end{tabular}
\end{table}

Per token per layer. Total reduction depends on sequence length and batch.

\subsection{Summary}

MLA reproduction:
\begin{enumerate}
\item Replaces separate K and V with a shared latent.
\item Initializes from concatenated K/V via SVD.
\item Normalizes singular values to preserve scale.
\item Handles four architectures with Conv1D transposition, nested bypass, and
  GQA expansion.
\item Uses LoRA for recovery.
\item Achieves \textbf{32$\times$} KV-cache reduction.
\end{enumerate}
\begin{lstlisting}

---
\end{lstlisting}